\documentclass[10pt,a4paper]{amsart}
\usepackage[margin=19mm]{geometry}
\usepackage{amsmath,amssymb,mathtools}
\usepackage[scr=rsfs,cal=euler]{mathalfa}
\usepackage{xfrac}
\usepackage[unicode,hidelinks,bookmarks=false]{hyperref}
\newcommand{\ie}{i.e.\ }
\newcommand{\cf}{cf.\ }
\newcommand{\resp}{respectively\ }
\newcommand{\Subsection}{\subsection}
\providecommand{\nobreakdash}{\nobreak\hbox{-}}
\setlength{\emergencystretch}{2em}
\multlinegap0pt
\usepackage{mathtools}
\usepackage{amssymb}
\usepackage{xcolor}
\usepackage{bm}
\newtheorem{proposition}{Proposition}
\newtheorem{lemma}{Lemma}
\def\B{\mathsf B}
\newcommand{\C}{{\mathbb C}}
\def\E{\mathsf E}
\def\Ext{\operatorname{Ext}}
\DeclareMathOperator{\Hom}{Hom}
\renewcommand{\P}{{\mathbb P}}
\def\W{\mathsf W}
\def\b{\mathsf b}
\def\cale{\mathcal E}
\def\calf{\mathcal F}
\def\e{\mathsf e}
\newcommand{\la}{\label}
\def\w{\mathsf w}
\newcommand{\wt}{\operatorname{wt}}
\title{Dimers and Beauville integrable systems}
\author{Terrence George, Giovanni Inchiostro}
\date{Source: arXiv:2207.09528v2 (CC BY 4.0)}
\begin{document}
\maketitle

\noindent\emph{Source and licence.} Dimers and Beauville integrable systems. Version arXiv:2207.09528v2 (CC BY 4.0), distributed by arXiv under the Creative Commons Attribution 4.0 International licence, http://creativecommons.org/licenses/by/4.0/, as recorded in the arXiv licence metadata for this exact version and shown on the abstract page https://arxiv.org/abs/2207.09528v2. Full licence notice: \texttt{licenses/arxiv-2207.09528-CC-BY-4.0.txt}.\par\smallskip

\noindent\textit{Editorial scope.} Counted unit: the article's lemma lemma (source0.tex lines 2637--2659). The article's complete original proof of it is delivered with it (source0.tex lines 2661--2783, one \texttt{proof} environment of 123 lines). No mathematics is written, repaired or re-derived: the statement and the proof are the article's own lines, in the article's own notation, and no retained line is substituted or re-flowed.  Retained uncounted as the source-local context the proof needs: lines 2139--2162; lines 2373--2403; lines 2453--2465.  Nothing was omitted from any retained span, so the package carries no omission record.  No retained line contains a citation, so the wrapper carries no bibliography and no bibliography entry.  No retained line contains a figure environment, an image-inclusion command or drawing code, so no image is delivered and the package declares no asset dependency.  Editorial formatting only, in addition: an AMS \texttt{amsart} preamble that restates the article's own declarations and macros used by the retained lines, namely the environments \texttt{proposition}, \texttt{lemma} on separate counters, together with the standard packages those lines need.  The retained environments print in this document as Proposition 1, Lemma 1 in order of appearance, whereas the article prints the same items with the numbering of its own full document.  The complete native source of the article as distributed in the arXiv e-print for this exact version is bundled unchanged as \texttt{sources/128447/source0.tex}.
\begin{proposition}\label{prop:repext}
There is an identification
\[
\Ext^\bullet_{\widetilde\P^2}(\widetilde{\mathcal L},\widetilde{\mathcal L})^{G}
\cong
H^\bullet\Bigl(\Bigl[
\begin{array}{c c c}
\C^{\B}\oplus \C^{\W} & \xrightarrow{\ d\ } & \C^{\E}\\[-.4ex]
{\scriptstyle 0} && {\scriptstyle 1}
\end{array}
\Bigr]\Bigr),
\]
where
\[
d(f,g)(\e)=\wt(\e)\bigl(f(\b)-g(\w)\bigr)
\qquad\text{for }\e=\b\w\in \E.
\]
In particular,
\[
\Ext^1_{\widetilde\P^2}(\widetilde{\mathcal L},\widetilde{\mathcal L})^{G}
\cong
\C^{\E}/\operatorname{im}(d).
\]
\end{proposition}

\begin{proposition}\label{prop:repext2}
There is an identification
\begin{equation}\label{eq:chain_ext-d}
\Ext^\bullet(\widetilde{\mathcal L},\widetilde{\mathcal L}(-\widetilde D))^G
\cong
H^\bullet\Bigl(\Bigl[
\begin{array}{c c c}
\C^{\E} & \xrightarrow{\ d^1\ } & \C^{\B}\oplus \C^{\W}\\[-.4ex]
{\scriptstyle 1} && {\scriptstyle 2}
\end{array}
\Bigr]\Bigr),
\end{equation}
where
\[
d^1(v)
=
\left(
\left(\sum_{\w:\b\w\in\E}\wt(\b\w)v(\b\w)\right)_{\b\in\B},
\left(\sum_{\b:\b\w\in\E}\wt(\b\w)v(\b\w)\right)_{\w\in\W}
\right).
\]
In particular,
\[
\Ext^1(\widetilde{\mathcal L},\widetilde{\mathcal L}(-\widetilde D))^G
\cong
\ker(d^1), \qquad
\Ext^2(\widetilde{\mathcal L},\widetilde{\mathcal L}(-\widetilde D))^G
\cong
(\C^{\B}\oplus \C^{\W})/\operatorname{im}(d^1).
\]
\end{proposition}

\begin{lemma}\label{lem:Ext2}
There is an identification
\begin{equation}\label{eq:Ext2-concrete}
\begin{aligned}
\Ext^\bullet(\widetilde{\mathcal L},\widetilde{\mathcal L}(-\widetilde D))^G
= H^\bullet\Bigl(\Bigl[ &
\underset{1}{H^2(\mathcal H om(\calf,\cale(-\widetilde D)))^G} \\
&\xrightarrow{\,d^{-1}\,}
\underset{2}{H^2(\mathcal H om(\cale,\cale(-\widetilde D))\oplus \mathcal H om(\calf,\calf(-\widetilde D)))^G}
\Bigr]\Bigr).
\end{aligned}
\end{equation}
\end{lemma}

\begin{lemma}\la{lem::cup1}
With the identifications in Proposition~\ref{prop:repext} and Proposition~\ref{prop:repext2}, the cup product map
\[
\cup:
\Ext^1(\widetilde{\mathcal L},\widetilde{\mathcal L})^G \otimes
\Ext^1(\widetilde{\mathcal L},\widetilde{\mathcal L}(-\widetilde D))^G
\rightarrow
\Ext^2(\widetilde{\mathcal L},\widetilde{\mathcal L}(-\widetilde D))^G
\]
is induced by the map
\[
\cup: \C^\E\otimes \C^\E \rightarrow \C^\B\oplus \C^\W
\]
given by
\[
u \cup v
=
\left(
\left(\sum_{\w:\b\w\in \E} u(\b\w)v(\b\w)\right)_{\b\in \B},
0
\right).
\]
\end{lemma}

\begin{proof}
By definition, the cup product is induced by the map
\[
\cup :
\check C^{p}\bigl(\widetilde{\mathcal U},\mathcal H om(K^\bullet,K^\bullet)^q\bigr)^G
\otimes
\check C^{p'}\bigl(\widetilde{\mathcal U},\mathcal H om(K^\bullet,K^\bullet(-\widetilde D))^{q'}\bigr)^G
\rightarrow
\check C^{p+p'}\bigl(\widetilde{\mathcal U},\mathcal H om(K^\bullet,K^\bullet(-\widetilde D))^{q+q'}\bigr)^G
\]
given by
\[
(\alpha\cup\beta)_{i_0\cdots i_{p+p'}}
=
(-1)^{q p'}
\beta_{i_p\cdots i_{p+p'}}\big|_{\widetilde U_{i_0\cdots i_{p+p'}}}
\circ
\alpha_{i_0\cdots i_p}\big|_{\widetilde U_{i_0\cdots i_{p+p'}}}.
\]

By Proposition~\ref{prop:repext}, an element \(u\in \C^\E\) determines an element of
\[
H^0(\mathcal H om(\cale,\calf))^G,
\]
represented by the \v Cech \(0\)-cocycle
\[
l=\bigl(u(\e)\,\widetilde x(\e)\bigr)_{\e\in\E}
\in \check C^0(\widetilde{\mathcal U},\mathcal H om(\cale,\calf))^G,
\]
where each \(u(\e)\,\widetilde x(\e)\) is viewed as a global section, hence as a \v Cech \(0\)-cocycle with identical components on the three open sets of \(\widetilde{\mathcal U}\). This class is represented by the $D$-cocycle
\[
(l,0,0)
\in
\check C^0(\widetilde{\mathcal U},\mathcal H om(\cale,\calf))^G
\oplus
\check C^1(\widetilde{\mathcal U},\mathcal H om(\cale,\cale)\oplus \mathcal H om(\calf,\calf))^G
\oplus
\check C^2(\widetilde{\mathcal U},\mathcal H om(\calf,\cale))^G,
\]
where \(\check d(l)=0\).

Similarly, let \(v\in \ker(d^1)\subset \C^\E\). A class in
\[
\Ext^1(\widetilde{\mathcal L},\widetilde{\mathcal L}(-\widetilde D))^G
\]
may be represented by a $D$-cocycle $(l',m',n')$ in
\[
\check C^0(\widetilde{\mathcal U},\mathcal H om(\cale,\calf(-\widetilde D)))^G
\oplus
\check C^1\!\bigl(\widetilde{\mathcal U},
\mathcal H om(\cale,\cale(-\widetilde D))
\oplus
\mathcal H om(\calf,\calf(-\widetilde D))\bigr)^G
\oplus
\check C^2(\widetilde{\mathcal U},\mathcal H om(\calf,\cale(-\widetilde D)))^G,
\]
satisfying
\begin{align*}
\check d(l')-d^0(m')=0,\qquad
\check d(m')+d^{-1}(n')=0,\qquad
\check d(n')=0.
\end{align*}
Here the \(\check C^2\)-component is
\[
n'
=
\frac{1}{\widetilde x_0 \widetilde x_1 \widetilde x_2}
\left( \frac{v(\e)}{\widetilde x(\e)} \right)_{\e\in\E}.
\]
The nonzero part of the cup product is
\[
(r,s)=(l,0,0)\cup(l',m',n'),
\]
where
\[
(r,s)\in
\check C^1(\widetilde{\mathcal U},\mathcal H om(\cale,\calf(-\widetilde D)))^G
\oplus
\check C^2\!\bigl(\widetilde{\mathcal U},
\mathcal H om(\cale,\cale(-\widetilde D))
\oplus
\mathcal H om(\calf,\calf(-\widetilde D))\bigr)^G
\]
is explicitly given by
\begin{align*}
r_{ij}
&=
-m_{ij}^{\prime,\calf}\circ l_i
\in
\Hom(\cale|_{\widetilde U_{ij}},\calf(-\widetilde D)|_{\widetilde U_{ij}})^G,\\[0.5ex]
s_{ijk}
&=
\bigl(n'_{ijk}\circ l_i,\ 0\bigr)
\in
\Hom(\cale|_{\widetilde U_{ijk}},\cale(-\widetilde D)|_{\widetilde U_{ijk}})^G
\oplus
\Hom(\calf|_{\widetilde U_{ijk}},\calf(-\widetilde D)|_{\widetilde U_{ijk}})^G.
\end{align*}

Under the identification of Lemma~\ref{lem:Ext2}, the class of the total cocycle \((r,s)\) is sent precisely to the class of \(s\) in
\[
H^2(\mathcal H om(\cale,\cale(-\widetilde D))\oplus \mathcal H om(\calf,\calf(-\widetilde D)))^G
\Big/
\operatorname{im}\bigl(d^{-1}\bigr).
\]
It remains to identify this class explicitly. The \((\b,\b)\)-entry of \(n'\circ l\) is
\[
\sum_{\w:\b\w\in \E}
\frac{u(\b\w)v(\b\w)}{\widetilde x_0\widetilde x_1\widetilde x_2}.
\]
Hence under the identification
\[
H^2(\mathcal H om(\cale,\cale(-\widetilde D))\oplus \mathcal H om(\calf,\calf(-\widetilde D)))^G
\cong \C^\B\oplus \C^\W,
\]
the class of \(s\) is
\[
\left(
\left(\sum_{\w:\b\w\in \E}u(\b\w)v(\b\w)\right)_{\b\in \B},
0
\right).
\]
\end{proof}

\end{document}
